The global financial crisis (GFC) of 2008 underscored the link between the financial sector and the broader economy, placing the concept of the \textit{financial cycle} at the forefront of academic and policy discussions. The financial cycle captures the self-reinforcing interactions between perceptions of value, risk-taking behaviour, and financing constraints, typically manifesting as periods of credit expansion followed by busts \citep{borio2014financial}. The expansion of the financial cycle, often reflected in rapid growth in credit, asset prices, and leverage, can lead to the accumulation of cyclical systemic risk and pose a threat to financial stability.

Accurately measuring the financial cycle is therefore central to macroprudential policy, particularly for the calibration of the countercyclical capital buffer (CCyB). The CCyB is designed to ensure that bank capital requirements reflect the prevailing level of cyclical systemic risk. During the expansionary phase of the financial cycle, when credit growth accelerates and asset prices increase, a gradual build-up of capital strengthens the resilience of the financial system. Conversely, during downturns, releasing the CCyB supports the flow of credit by allowing banks to absorb losses while sustaining lending, thereby mitigating the amplification of adverse shocks to the real economy. In practice, however, cyclical systemic risk is difficult to assess in real time, and macroprudential policy tools such as the CCyB are subject to implementation lags. For this reason, many jurisdictions have adopted a positive neutral CCyB rate even during periods of moderate systemic risk, ensuring that releasable capital is available when financial conditions deteriorate.

Against this backdrop, this paper contributes to the literature on financial cycle measurement by developing a simple methodology to construct country-specific indicators of the financial cycle. The proposed approach captures the common dynamics across key variables reflecting credit developments, asset price fluctuations, leverage, and risk-taking behaviour in the financial system. Specifically, we employ a dynamic factor model (DFM) to extract a latent factor summarising the co-movement of financial variables for a sample of European countries. The use of a common factor is grounded in theoretical frameworks that emphasise self-reinforcing mechanisms in the financial sector, including the financial accelerator \citep{bernanke1999financial}, collateral constraints \citep{kiyotaki1997credit}, and correlated risk-taking by financial intermediaries \citep{he2013intermediary}.

Several studies have applied factor models to analyse the financial cycle, including \cite{jorg2016analyzing, menden2017dissecting, miranda2020us}. However, existing contributions predominantly focus on global financial cycle indicators and their relationship with macroeconomic conditions or monetary policy. In contrast, we develop country-specific measures that are directly relevant for macroprudential policy and explicitly link the financial cycle to cyclical systemic risk and systemic banking crises, with a particular emphasis on early-warning properties.

The paper makes three main contributions to this literature. First, the proposed indicator provides an accurate and coherent description of the evolution of the financial cycle across countries, consistent with prior expectations from the literature and historical experience. The estimated cycle identifies three prominent peaks (in the early 2000s, in the run-up to the Global Financial Crisis, and following the COVID-19 shock) as well as the prolonged trough associated with post-GFC deleveraging and the renewed contraction from 2022 onwards as monetary policy tightened globally. These patterns align closely with known episodes of financial expansion and stress, lending credibility to the indicator as a measure of cyclical systemic risk. The indicator can further be decomposed into contributions from individual variables, providing an interpretable account of the main drivers of financial cycle dynamics at each phase. By revealing which components (credit, asset prices, indebtedness, or financial conditions) are the primary sources of cyclical risk at any given point, the decomposition enhances the transparency of the indicator and allows macroprudential authorities to identify the nature and origin of financial imbalances and to communicate their assessments in a structured manner.

Second, and most distinctively, the paper introduces a statistically grounded framework for identifying periods of neutral cyclical systemic risk. By drawing on the asymptotic distribution of the principal components estimator, the approach yields time-varying confidence bands that distinguish elevated, neutral, and subdued financial cycle regimes without relying on ad hoc thresholds or judgement. This addresses a practical gap for macroprudential authorities: existing approaches offer no principled, data-driven criterion for determining when cyclical risk is neither building up nor materialising, which is precisely the condition underpinning the positive neutral CCyB. The proposed framework fills this gap by providing a transparent, statistically grounded benchmark for regime classification, offering a rigorous and operational basis for CCyB calibration that is anchored in the data rather than in discretionary assessments.

Third, and of direct relevance for policy applications, the indicator is evaluated in a pseudo real-time, out-of-sample exercise designed to replicate the information constraints faced by policymakers in practice. The results point to good early-warning properties. The estimated one-year-ahead probability of a systemic banking crisis approximately doubles in the two years preceding the Global Financial Crisis, rising from around 10\% in early 2006 to approximately 20\% by 2007, a signal generated entirely from information available in real time and without any knowledge of the forthcoming crisis. In-sample results further confirm that the DFM-based indicator outperforms the Basel gap, the standard regulatory benchmark, in terms of discriminatory power, as measured by the area under the receiver operating characteristic curve. Together, these findings strengthen the case for using the proposed indicator as an operational tool to monitor cyclical systemic risks. Throughout, the methodology is deliberately simple and transparent, making it easy to implement in practice and suitable as a benchmark for policymakers.

This paper contributes to a broad literature on the measurement of the financial cycle, where a range of approaches have been proposed. A widely used benchmark is the credit-to-GDP gap developed by the Bank for International Settlements (BIS). While effective in providing early-warning signals of financial crises, its reliance on the Hodrick--Prescott (HP) filter introduces biases in trend estimation. In particular, the expansion phase of the financial cycle tends to be excessively incorporated into the estimated trend, leading to prolonged periods with negative values of the credit-to-GDP gap following episodes of systemic risk build-up. These limitations have motivated alternative methodologies. For example, \cite{drehmann2012characterising} employ turning-point analysis to identify local extrema, while frequency-based filters reduce reliance on assumptions about the characteristics of the financial cycle but introduce arbitrariness in filter selection. Composite indicators, such as \cite{lang2019anticipating}, are designed as early-warning tools for systemic banking crises, whereas \cite{schuler2020financial} employ spectral methods to capture co-movement across variables without imposing strong parametric assumptions.

The remainder of the paper is structured as follows. Section~\ref{Data} presents the data and the rationale for variable selection. Section~\ref{Method} outlines the dynamic factor model methodology. Section~\ref{Results} presents the empirical results, and Section~\ref{Conclusion} concludes.
